\section{Further research questions with concrete mathematical targets}
\label{future:sec}

The following are proposals for further work, not consequences claimed
in the preceding theorems. They prioritize exact structural statements
and finite evaluations. A question asking for a classification allows
a counterexample with a proved exceptional set as a useful answer.

\subsection{Reconstruction beyond order ten}

\begin{question}[All unequal even orders]
For every $1\le a<b$, is the first-$2b$ centered even-moment prefix
of the $(2a,2b)$ tail product nonsingular in its formal source
coordinates? The present article completes $a=5$; the project now
has all $a\le5$. A decisive next result would be a structural
noncancellation theorem for the fixed-size Bernoulli determinant,
or a proved description of any exceptional pairs.
\end{question}

\begin{question}[A uniform overlap principle]
Can one relate the precision of a doubled-Bernoulli congruence to the
constructive real-variable threshold for arbitrary fixed $a$?
At $a=5$ the sole exceptional class has least positive representative
$3010$, exceeding the threshold $1920$. Higher precision alone does
not prove nonvanishing: a compatible $2$-adic root may persist. The
target is a theorem ensuring coverage by two complementary arguments,
not a table of increasingly long nonzero prefixes.
\end{question}

\begin{question}[The full gap-six family]
Give a bordered-Hankel reduction that retains the moving $B_0$ and
$B_1$ endpoint terms for $b=a+3$. The stable fixed-order contraction
requires $b\ge2a$ and therefore does not cover this family uniformly.
A formula with all boundary terms retained could settle the global
gap-six conjecture without extrapolating the finitely checked pairs.
\end{question}

\begin{question}[Compact inverse rows]
Find a product, determinant of uniformly smaller size, or short
recurrence for the polynomials $Q_{b,k}$ of
Corollary~\ref{ten:identity}. Their exact rational computability is
proved, but the displayed $(10,12)$ example shows that naive expanded
coefficients grow rapidly. A compact expression should also identify
the subtraction polynomial directly and preserve ordinary convergence.
\end{question}

\subsection{Beyond the bilinear Gamma normalization}

\begin{question}[Every bilinear contact hyperplane]
For the two independent orders in Theorem~\ref{bil:master}, compute
the constant terms of $C_N$ and $(-1)^ND_N$ on the further resonant
hyperplanes $p+q=-2m$. The first symmetric contact is $1/2$ at
$p+q=0$. Seek closed bivariate contact polynomials, every independent
$p,q$ jet, and the exact change under translating the cutoff.
The diagonal answer in Section~\ref{cut:sec} is a mandatory check;
it does not determine all transverse derivatives.
\end{question}

\begin{question}[Independent outer weights and ordinary Hurwitz tails]
Construct a finite-cutoff identity for
$\zeta(p,n+a)\zeta(q,n+b)(n+c)^{-w}$ with independent
$p,q,w,a,b,c$, retaining an explicit remainder or a precisely defined
finite collection of additional functions. Determine its moving
polar hyperplanes and residues before differentiating on them.
The Gamma-weighted alternating identities do not settle this ordinary
Hurwitz problem; a valid result must state the size and analytic
meaning of any newly required coordinate collection.
\end{question}

\begin{question}[More colors and higher products]
For tails with a primitive $d$th-root twist, can a product of Gamma
quotients convert the $d$-step recurrence into a finite bilinear
endpoint law? For three tail factors, identify the smallest adjacent
product ansatz whose linear and quadratic residual terms can all
cancel. A useful target is a divisor-orbit criterion extending
Theorem~\ref{rat:classification}, together with one ordinary convergent
three-factor evaluation that genuinely lies outside its bilinear case.
\end{question}

\begin{question}[The remaining logarithmic Gamma generator]
The logarithmic square sum in Corollary~\ref{cut:Gamma} is explicit,
but the earlier twisted logarithmic-Gamma correction generator has
an independent outer resolvent. Determine whether it belongs to a
finite stated collection of classical or multiple-Gamma functions,
or prove a functional obstruction in a declared difference-equation
class. A failure of numerical reduction is not such an obstruction.
\end{question}

\subsection{Harmonic pencils, spectral coefficients, and polylogarithms}

\begin{question}[Complex weights and sharp cutoff exceptions]
Does generic degree--rank equality extend to every complex rational
weight? If not, classify the smallest counterexamples using the
kernel-compression identity \eqref{deg:highest}. For real weights,
determine whether the bound $r(r-\rho)$ on late deficient cutoffs is
sharp in the rational-moment class, and how it improves for special
polynomial kernel spaces. The real proof uses definite increments;
it cannot be transferred to a complex bilinear Gram form by wording
alone.
\end{question}

\begin{question}[Higher-rank pencils and their realizations]
Classify rational residue data whose entire harmonic pencil has
generic rank two, then rank $\rho$. The family in
Theorem~\ref{hier:rank} realizes every rank, but does not classify all
pencils: the maximal-rank kernel need not be a common kernel.
Seek exact finite residue conditions and identify which such pencils
come from positive rational weights. This would complement the
existing rank-one classification.
\end{question}

\begin{question}[All Laurent coefficients in compact form]
For the hierarchy $Z_{r,\rho}(s,a)$, find short formulas or recurrences
for the lower normalizations $G_{L,h}^{(m)}(a)$ at $m\ge2$, and for
the complete principal and finite parts at every pole. The leading
coefficient at $m=1$ is a Gamma product. At higher orders, confluent
deformations of the mixed factorial moments may reduce determinant
growth. Any claimed formula should reproduce both finite parts in
\eqref{hier:r2-s1}--\eqref{hier:r2-s2} and distinguish them from
hard-cutoff constants.
\end{question}

\begin{question}[Complementary cyclotomic generators and exact depth]
Construct the full complementary two-resolvent generator requested
in the alternating report, retaining the evaluated coefficient
\eqref{audit:coefficient-value} as a boundary check. In parallel,
determine minimal depth for concrete harmonic-determinant generating
functions inside a specified shuffle, differential, or symbol
calculus. Formal harmonic degree predicts a useful representation
bound; it does not by itself prove minimal polylogarithmic depth.
The Gaussian $S_6$ and revised $S_8$ assertions remain separate
conjectures requiring exact reductions or counterexamples.
\end{question}

\subsection{A practical order of attack}

The next finite target is smaller order twelve: derive its exact
fixed-size numerator, determine a sufficiently precise Bernoulli
congruence, and compare the exceptional classes with the real-variable
threshold. The next analytic target is the first off-diagonal contact
hyperplane $p+q=-2$, where the new bilinear endpoint contains enough
independent parameters to reveal which diagonal parity cancellations
survive. The next determinant target is a rank-two residue
classification, tested against both the positive shorter-stencil
family and a pencil without a common kernel.

For formal verification, the finite bilinear recurrence, rational
root-orbit criterion, and symmetric-pencil compression lemma are
small natural starting points. The order-ten congruence can then be
checked in a polynomial ring modulo $2^{17}$ from the finite Newton
matrix. Such formalization would strengthen the trust chain without
altering the claimed mathematics or replacing a proof by decimal
agreement.
